\chapter{Implementacja mikroserwisów}
\label{chap:implementacja-mikroserwisow}

Rozdział ten przedstawia implementację trzech mikroserwisów systemu w~języku Go. Omówiono wspólną strukturę
i~wzorce dzielone przez wszystkie usługi, a~następnie logikę biznesową każdej z~nich, warstwę zdarzeń, mechanizm
uwierzytelniania, obserwowalność, konteneryzację oraz testy. Prezentowane listingi pochodzą bezpośrednio z~kodu
źródłowego systemu.

\section{Struktura projektu i wspólne wzorce serwisów}
\label{sec:struktura-projektu}

Wszystkie trzy serwisy zbudowano według tego samego szablonu, opartego na architekturze heksagonalnej
(\emph{ports and adapters}). Logika domenowa stanowi rdzeń niezależny od technologii. Komunikacja ze światem
zewnętrznym przebiega przez HTTP. Baza danych, broker komunikatów odbywa się poprzez wymienne adaptery, sprzęgnięte
z~rdzeniem wyłącznie interfejsami. Strukturę katalogów serwisu zamówień przedstawia listing~\ref{lst:struktura-serwisu}.

\begin{lstlisting}[caption={Struktura katalogów serwisu zamówień},label={lst:struktura-serwisu}]
order-service/
|-- main.go               # punkt wejścia, złożenie zależności
|-- Dockerfile
`-- internal/
    |-- auth/             # weryfikacja JWT (Cognito)
    |-- config/           # konfiguracja ze zmiennych środowiskowych
    |-- db/               # pula połączeń PostgreSQL, migracje
    |-- metrics/          # definicje metryk Prometheus
    |-- producer/         # producent zdarzeń Kafka
    |-- server/           # serwery HTTP: API i techniczny
    `-- order/            # kontekst domenowy
        |-- domain/       # encje, maszyna stanów, niezmienniki
        |-- service/      # przypadki użycia (logika aplikacyjna)
        |-- repository/   # adapter PostgreSQL
        |-- events/       # definicje zdarzeń domenowych
        `-- api/          # handlery HTTP, DTO, walidacja wejścia
\end{lstlisting}

Zależności składane są jawnie w~funkcji \code{main}, bez wstrzykiwania zależności. Granice pomiędzy
warstwami wyznaczają małe interfejsy definiowane po stronie konsumenta, co jest idiomatycznym wzorcem języka Go.
Tworzy to pakiety bez zewnętrznych zależności co pozwala testować logikę biznesową w odpowiedniej izolacji.

Każdy serwis uruchamia dwa niezależne serwery HTTP na oddzielnych portach. Serwer \textbf{API} obsługuje ruch biznesowy. Serwer
\textbf{techniczny} udostępnia punkty końcowe eksploatacyjne: \code{/ping} (kontrola zdrowia), 
\code{/metrics} (metryki Prometheus), \code{/version} oraz \code{/debug/pprof} (profilowanie procesu).
Rozdzielenie portów gwarantuje, że dane eksploatacyjne nigdy nie są osiągalne ścieżką publiczną, nawet przy błędnej konfiguracji routingu.

Konfiguracja serwisów realizuje wzorzec konfiguracji zewnętrznej (sekcja~\ref{sec:wzorce-projektowe}). Wszystkie
parametry wczytywane są ze zmiennych środowiskowych do struktur Go.

\section{Serwis zamówień - logika biznesowa, maszyna stanów i idempotencja}
\label{sec:order-service}

Serwis zamówień realizuje utworzenie zamówienia oraz pobranie i~listowanie zamówień klienta.
Tworząc zamówienie, warstwa domenowa egzekwuje niezmienniki: niepustą listę pozycji,
dodatnie ilości, poprawny kod waluty oraz zgodność kwoty całkowitej z~sumą pozycji.

Centralnym elementem domeny jest maszyna stanów zamówienia. Dozwolone przejścia zakodowano jako strukturę danych.
Mapę stanów osiągalnych z~każdego stanu, dzięki czemu cykl życia zamówienia opisany jest w~jednym miejscu kodu 
i~bezpośrednio odpowiada diagramowi z~projektu (sekcja~\ref{sec:model-domenowy}). Implementację przedstawia 
listing~\ref{lst:maszyna-stanow}. Stany terminalne nie posiadają wpisu w~mapie, więc żadne przejście z~nich nie jest możliwe.

\begin{lstlisting}[caption={Implementacja maszyny stanów zamówienia},label={lst:maszyna-stanow},language=Go]
// allowedTransitions defines the order lifecycle.
//
//  NEW -> VALIDATED -> PAYMENT_PENDING -> PAID -> PROCESSING -> SHIPPED
//  plus CANCELLED / FAILED reachable from any non-terminal state.
var allowedTransitions = map[OrderStatus]map[OrderStatus]struct{}{
    OrderStatusNew:            {OrderStatusValidated: {}, OrderStatusCancelled: {}, OrderStatusFailed: {}},
    OrderStatusValidated:      {OrderStatusPaymentPending: {}, OrderStatusCancelled: {}, OrderStatusFailed: {}},
    OrderStatusPaymentPending: {OrderStatusPaid: {}, OrderStatusCancelled: {}, OrderStatusFailed: {}},
    OrderStatusPaid:           {OrderStatusProcessing: {}, OrderStatusCancelled: {}, OrderStatusFailed: {}},
    OrderStatusProcessing:     {OrderStatusShipped: {}, OrderStatusCancelled: {}, OrderStatusFailed: {}},
}

func CanTransition(from, to OrderStatus) bool {
    if from == to {
        return false
    }
    next, ok := allowedTransitions[from]
    if !ok {
        return false
    }
    _, ok = next[to]
    return ok
}
\end{lstlisting}

Każde wykonane przejście dopisywane jest do historii zamówienia (stan poprzedni, nowy, znacznik czasu, opcjonalny
powód), utrwalanej w~kolumnie \code{JSONB}. Daje to pełny ślad cyklu życia każdego zamówienia.

Idempotencję operacji utworzenia zamówienia zrealizowano na poziomie bazy danych, w~sposób odporny na wyścigi
współbieżnych żądań (listing~\ref{lst:idempotencja}). Klucz z~nagłówka \code{Idempotency-Key} zapisywany jest
instrukcją \code{INSERT} z~klauzulą \code{ON CONFLICT}. Jeśli klucz został już zajęty, zapytanie zamiast
zakończyć się błędem zwraca identyfikator zamówienia utworzonego pierwotnie, a~serwis odpowiada klientowi tym
właśnie zamówieniem. Rozstrzygnięcie kolizji następuje zatem atomowo w~jednym zapytaniu, bez okna czasowego
pomiędzy sprawdzeniem a~zapisem klucza.

\begin{lstlisting}[caption={Atomowe zajęcie klucza idempotencji},label={lst:idempotencja},language=Go]
func (r *Postgres) SaveIdempotencyKey(ctx context.Context, key, orderID string) (string, error) {
    var stored string
    err := r.pool.QueryRow(ctx, `
        INSERT INTO idempotency_keys (key, order_id) VALUES ($1, $2)
        ON CONFLICT (key) DO UPDATE SET key = idempotency_keys.key
        RETURNING order_id
    `, key, orderID).Scan(&stored)
    if err != nil {
        return "", fmt.Errorf("save idempotency key: %w", err)
    }
    if stored != orderID {
        return stored, domain.ErrIdempotencyExists
    }
    return "", nil
}
\end{lstlisting}

Po utrwaleniu zamówienia serwis publikuje zdarzenie \code{order.created}.

\section{Serwis produktów - katalog i rezerwacja stanów magazynowych}
\label{sec:product-service}

Serwis produktów pełni dwie role. Synchronicznie obsługuje katalog (listowanie z~filtrowaniem po kategorii,
przedziale cenowym i~dostępności, obsługę zdjęć przechowywanych w~kolumnie \code{BYTEA}).
Oraz asynchronicznie uczestniczy w~walidacji zamówień jako konsument zdarzeń.

Rezerwacja stanów magazynowych realizowana jest operacją wsadową obejmującą wszystkie pozycje zamówienia w~jednej
transakcji bazodanowej. Rezerwacja jest więc niepodzielna. Nie istnieje stan, w~którym część pozycji zamówienia
została zarezerwowana, a~część nie. Samo przesunięcie ilości pomiędzy pulami wykonywane jest warunkową instrukcją
\code{UPDATE}. Warunek (\code{qty >= \$2}) gwarantuje poprawność również przy współbieżnej rezerwacji tego
samego produktu przez dwa zamówienia. Przegrywające zapytanie nie zmodyfikuje żadnego wiersza, co serwis
interpretuje jako niewystarczający stan magazynowy.

Po otrzymaniu zdarzenia \code{order.created} konsument dekoduje je, waliduje i~podejmuje próbę rezerwacji.
W~zależności od wyniku publikuje \code{order.validated} albo \code{order.reservation\_failed}.

Rozwiązaniem przejściowym, jest sposób domknięcia pętli walidacji. Docelowo order-service powinien 
konsumować temat \code{product.events} i~sam zmieniać status zamówienia. W~bieżącej
wersji systemu konsument ten nie istnieje. Product-service aktualizuje więc status zamówienia
(\code{VALIDATED} lub \code{CANCELLED}) bezpośrednio w~tabeli zamówień, równolegle z~publikacją zdarzenia.
Kompromis ten narusza zasadę wyłącznej własności danych (sekcja~\ref{sec:charakterystyka-mikroserwisow})
i~został świadomie przyjęty jako etap pośredni. Jego usunięcie wskazano jako kierunek rozwoju
w~podsumowaniu pracy.

\section{Serwis płatności - transakcyjna obsługa procesu płatności}
\label{sec:payment-service}

Serwis płatności obsługuje utworzenie żądania płatności dla zamówienia (zwracając klientowi adres URL płatności)
oraz jej finalizację. Finalizacja jest najbardziej krytyczną operacją systemu. Musi jednocześnie i~niepodzielnie
przenieść zarezerwowane ilości wszystkich pozycji zamówienia do puli opłaconych, oznaczyć zamówienie jako
\code{PAID} oraz płatność jako \code{SUCCEEDED}. Zrealizowano ją jako pojedynczą transakcję bazodanową
z~blokadami wierszy.

Transakcja rozpoczyna się pobraniem płatności i~zamówienia z~klauzulą \code{FOR UPDATE}, co wyklucza równoległą
finalizację tej samej płatności. Następnie weryfikowane są warunki: przynależność płatności i~zamówienia do
uwierzytelnionego klienta, poprawność pozycji oraz to, czy zamówienie znajduje się w~stanie dopuszczającym
zapłatę. Przy czym ponowna finalizacja płatności już zakończonej nie jest błędem, lecz zwraca jej bieżący stan
(operacja jest idempotentna). Dla każdej pozycji zamówienia wykonywane jest warunkowe przesunięcie ilości
(listing~\ref{lst:platnosc-update}). Niepowodzenie któregokolwiek przesunięcia wycofuje całą transakcję.

\begin{lstlisting}[caption={Warunkowe przesunięcie ilości w transakcji płatności},label={lst:platnosc-update},language=Go]
tag, err := tx.Exec(ctx, `
    UPDATE products
    SET reserved_qty = reserved_qty - $2,
        paid_qty = paid_qty + $2
    WHERE id = $1 AND reserved_qty >= $2
`, item.ProductID, item.Quantity)
\end{lstlisting}

Transakcja ta obejmuje tabele trzech kontekstów domenowych (płatności, zamówień i~produktów), co stanowi opisany
w~sekcji~\ref{sec:architektura-systemu} świadomy kompromis względem wzorca \emph{database-per-service}.
Współdzielenie jednej instancji baz gwarantujący silną spójność najbardziej wrażliwej operacji
systemu bez implementacji rozproszonej sagi.

\section{Publikacja i konsumpcja zdarzeń Kafka}
\label{sec:implementacja-kafka}

Warstwę zdarzeń zrealizowano biblioteką \code{franz-go}. Producent serializuje zdarzenia do JSON i~publikuje je
do skonfigurowanego tematu w~trybie nieblokującym. Publikacja nie wydłuża obsługi żądania HTTP, ponieważ rekordy
trafiają do wewnętrznego bufora klienta i~są wysyłane wsadowo w~tle. Listing~\ref{lst:kafka-producer} przedstawia
konfigurację producenta w~order-service oraz metodę publikującą zdarzenie. Krótki czas zwłoki
(\code{ProducerLinger}) pozwala grupować rekordy w~większe wsady, ograniczenie bufora
(\code{MaxBufferedRecords}) chroni proces przed nieograniczonym wzrostem pamięci przy niedostępności brokera,
a~partycjonowanie po kluczu zachowuje kolejność zdarzeń dotyczących tego samego zamówienia.

\begin{lstlisting}[caption={Konfiguracja i nieblokująca publikacja zdarzeń w order-service},label={lst:kafka-producer},language=Go]
func kafkaOptionsNotifier(topic, urls string) []kgo.Opt {
    return []kgo.Opt{
        kgo.SeedBrokers(strings.Split(urls, ",")...),
        kgo.DefaultProduceTopic(topic),
        kgo.ProducerLinger(time.Duration(50) * time.Millisecond),
        kgo.MaxBufferedRecords(20_000),
        kgo.RecordPartitioner(kgo.StickyKeyPartitioner(nil)),
    }
}

func (n *EventNotifier) Notify(log any) error {
    if !n.enabled {
        return nil
    }
    b, err := json.Marshal(log)
    if err != nil {
        return fmt.Errorf("unable to marshal log: %w", err)
    }
    n.producer.Produce(context.Background(), kgo.SliceRecord(b), nil)

    return nil
}
\end{lstlisting}

Konsument w~product-service działa w~ramach grupy konsumentów \code{product-service}, dzięki czemu pozycja odczytu
tematu \code{order.events} jest trwała. Po restarcie serwis wznawia przetwarzanie od ostatniego zatwierdzonego
przesunięcia, nie gubiąc zdarzeń opublikowanych w~czasie jego niedostępności. Pętla odczytu
(listing~\ref{lst:kafka-consumer}) pobiera wsady rekordów i~przekazuje każdy z~nich do funkcji obsługi
przekazanej przy konstrukcji konsumenta. Błąd obsługi pojedynczego rekordu jest logowany, lecz nie zatrzymuje
przetwarzania kolejnych zdarzeń.

\begin{lstlisting}[caption={Pętla konsumenta zdarzeń w product-service},label={lst:kafka-consumer},language=Go]
opts := []kgo.Opt{
    kgo.SeedBrokers(strings.Split(cfg.URL, ",")...),
    kgo.ConsumerGroup(cfg.Group),
    kgo.ConsumeTopics(cfg.Topic),
}

func (c *Consumer) Run(ctx context.Context) error {
    ...
    defer c.client.Close()

    slog.Info("kafka consumer started")
    for {
        if ctx.Err() != nil {
            return nil
        }

        fetches := c.client.PollFetches(ctx)
        if fetches.IsClientClosed() {
            return nil
        }
        fetches.EachError(func(topic string, partition int32, err error) {
            if errors.Is(err, context.Canceled) {
                return
            }
            slog.Error("kafka fetch error", "topic", topic,
                "partition", partition, "err", err)
        })

        fetches.EachRecord(func(r *kgo.Record) {
            if err := c.handle(ctx, r.Value); err != nil {
                slog.Error("handle kafka record failed", "topic", r.Topic,
                    "offset", r.Offset, "err", err)
            }
        })
    }
}
\end{lstlisting}

\section{Walidacja tokenów JWT i integracja z Cognito}
\label{sec:implementacja-auth}

Uwierzytelnianie żądań zrealizowano jako middleware HTTP, włączany przed handlery wszystkich chronionych punktów
końcowych. Middleware wydobywa token z~nagłówka \code{Authorization}, przekazuje go do weryfikatora, a~po pomyślnej
weryfikacji umieszcza tożsamość klienta w~kontekście żądania. Skąd odczytują ją handlery, zawężając operacje 
do zasobów uwierzytelnionego użytkownika.

Produkcyjny weryfikator (listing~\ref{lst:cognito-verify}) opiera się na bibliotece \code{jwx}. Klucze publiczne
puli użytkowników pobierane są ze~standardowego adresu JWKS i~buforowane lokalnie z~okresowym odświeżaniem
(od~15~minut do~24~godzin).

\begin{lstlisting}[caption={Weryfikacja tokenu JWT wobec kluczy Cognito},label={lst:cognito-verify},language=Go]
func (v *CognitoVerifier) Verify(_ context.Context, raw string) (jwt.Token, error) {
    tok, err := jwt.Parse(
        []byte(raw),
        jwt.WithKeySet(v.keySet),
        jwt.WithValidate(true),
        jwt.WithIssuer(v.issuer),
        jwt.WithAcceptableSkew(30*time.Second),
    )
    if err != nil {
        return nil, fmt.Errorf("token verify: %w", err)
    }

    var tokenUse string
    if err := tok.Get("token_use", &tokenUse); err != nil {
        return nil, fmt.Errorf("token_use claim missing: %w", err)
    }
    if _, ok := v.allowedTokenUse[tokenUse]; !ok {
        return nil, fmt.Errorf("token_use %q not allowed", tokenUse)
    }
    ...
}
\end{lstlisting}

Na potrzeby pracy lokalnej przewidziano tryb deweloperski (\code{COGNITO\_ENABLED=false}), w~którym tożsamość
klienta przekazywana jest jawnie w~nagłówku żądania. Tryb ten nie jest dostępny w~konfiguracji produkcyjnej.

\section{Obserwowalność serwisów}
\label{sec:obserwowalnosc-serwisow}

Każdy serwis publikuje metryki Prometheus w~przestrzeni nazw odpowiadającej usłudze (prefiksy \code{order\_},
\code{product\_}, \code{payment\_}). Warstwa HTTP instrumentowana jest middleware'em zliczającym żądania, mierzącym histogram czasów
obsługi oraz liczbę żądań w~toku. Obok metryk technicznych serwisy publikują metryki domenowe, które pozwalają obserwować system
w~kategoriach biznesowych, a~nie wyłącznie infrastrukturalnych.

Logowanie zrealizowano pakietem standardowym \code{slog} z~formatem JSON i~konfigurowalnym poziomem. Serwer
techniczny każdej usługi udostępnia ponadto kontrolę zdrowia \code{/ping}, informację o~wersji \code{/version}
oraz profilowanie \code{pprof}, pozwalające diagnozować zużycie procesora i~pamięci działającego procesu.

\section{Konteneryzacja serwisów}
\label{sec:konteneryzacja-serwisow}

Obrazy kontenerów budowane są wieloetapowo (listing~\ref{lst:dockerfile}). Etap budowania, oparty na obrazie
\code{golang:alpine}, kompiluje statycznie linkowany plik wykonywalny. Osobne kopiowanie plików
\code{go.mod} i~\code{go.sum} oraz montowane pamięci podręczne modułów i~kompilatora pozwalają przy niezmienionych
zależnościach pominąć ich ponowne pobieranie. Obraz docelowy to \code{distroless/static}. Czyli obraz pozbawiony
powłoki, menedżera pakietów i~jakichkolwiek narzędzi systemowych. Zawierający wyłącznie skompilowany plik binarny
i~certyfikaty TLS. Uruchamiany jest jako użytkownik nieuprzywilejowany (\code{nonroot}). Minimalizuje to zarówno
rozmiar obrazu, jak i~powierzchnię ataku.

\begin{lstlisting}[caption={Wieloetapowy Dockerfile serwisu zamówień},label={lst:dockerfile},language=bash]
FROM golang:1.26.3-alpine AS build
WORKDIR /src

COPY go.mod go.sum ./
RUN --mount=type=cache,target=/go/pkg/mod \
    go mod download

COPY . .
RUN --mount=type=cache,target=/go/pkg/mod \
    --mount=type=cache,target=/root/.cache/go-build \
    CGO_ENABLED=0 GOOS=linux go build \
    -trimpath -ldflags="-s -w" \
    -o /out/order-service .

FROM gcr.io/distroless/static-debian12:nonroot
COPY --from=build /out/order-service /order-service
EXPOSE 9002 9090
USER nonroot
ENTRYPOINT ["/order-service"]
\end{lstlisting}

\section{Testy jednostkowe i integracyjne}
\label{sec:testy-backend}

Testy serwisów napisano w~standardowym frameworku testowym Go. W~konwencji testów tabelarycznych. Pokrywają one
dwa poziomy. Poziom domenowy weryfikuje logikę biznesową w~pełnej izolacji.

Testy uruchamiane są z~włączonym detektorem wyścigów (\code{go test -race}) oraz losową kolejnością wykonywania
(\code{-shuffle=on}). Pozwala to wykrywać odpowiednio błędy współbieżności i~ukryte zależności pomiędzy testami.
Te same polecenia wykonywane są automatycznie przy każdej zmianie kodu w~potoku ciągłej integracji. 
Testy obciążeniowe systemu wdrożonego w~środowisku docelowym omówiono odrębnie w~rozdziale~\ref{chap:testy-i-analiza}.